\documentclass[a4paper]{article}
\usepackage[pdftex]{graphicx}
\usepackage[utf8]{inputenc}
\usepackage{enumerate}
\usepackage{mdwlist}
\usepackage{amssymb}
\usepackage{icomma}
\usepackage{siunitx}
\sisetup{locale=DE}
\usepackage{href-ul}
\hypersetup{
	colorlinks=true,
	linkcolor=blue,
	urlcolor=blue}
\usepackage{geometry}
\geometry{a4paper, top=15mm, left=15mm, right=15mm, bottom=15mm,
	headsep=10mm, footskip=12mm}

\begin{document}
	{\bf Free fall and vertical throw}
	
	\begin{enumerate}[1.]
		{\bf \item According to legend, Sir Isaac Newton (1643–1727) } was lying under an apple tree when he watched an apple fall and jumped up from it
		his power law concluded. How much reaction time would he have had to catch the apple if the height of the fall had been $h=\SI{2.4}{\meter}$?
		How fast would the apple have grown then?
		
		{\bf \item A parcel drone comes flying and delivers} a package for Ms. Eilig, who is already standing in front of her house. What is the minimum flight altitude of the drone from which it drops the package if Ms. Eilig is $\SI{1.70}{\meter}$ tall and $\SI{0.5}{\second}$ reaction time needed to catch the package? How much kinetic energy does the $\SI{1.5}{\kilogram}$ package have when caught?
		
		{\bf \item The water from the fountain in the Künettegraben} flows out at a speed of $\SI{9,0}{\meter\over\second}$. How high is the fountain?
		How long is a drop of water in the jet in the air before it comes back down?
		
		{\bf \item TEXUS -- Technical experiments under weightlessness} For this purpose, rockets are fired vertically upwards in Kiruna/northern Sweden. During ballistic flight, the payload is $\SI{6.0}{\minute}$ almost in zero gravity.
		\begin{enumerate}[(a)]
			\item What flight altitude is theoretically required?
			\item In fact, the flight altitude reached is $\SI{250}{\kilo\meter}$. Give two reasons for this!
		\end{enumerate}
		
		{\bf \item Unfortunately the penalty hit the crossbar and the ball bounced vertically upwards.} The spidercam is directly $\SI{3.0}{\meter}$ above it and is filming the whole thing.
		\begin{enumerate}[(a)]
			\item If the camera is hit if the ball was fast before the rebound $v=\SI{11}{\meter\over\second}$ and through the rebound
			Loses 50% of its kinetic energy? (Try without calculating the mass -- it's possible!)
			\suspend{enumerate}
			With a second penalty, the ball is shot again with the same force and onto the crossbar, only this time it bounces downwards.
			\resume{enumerate}
			\item What is the kinetic energy when it hits the goalkeeper lying on the ground if the crossbar is $\SI{2.4}{\meter}$ above the ground and the mass of the ball is $m=\SI{0.42 }{\kilogram}$ is ?
			\item What speed does this correspond to? Calculate again without mass!
		\end{enumerate}
		
		
		{\bf \item Lunerald stands on the three-meter board.} His center of gravity is initially at a height of $h=\SI{4,0}{\meter}$ above the water and his mass is $m= \SI{75}{\kilogram}$. When jumping off, the board stretches with the spring hardness $D=\SI{7.5}{\kilo\newton\over\meter}$
		by $s=\SI{0.40}{\meter}$
		\begin{enumerate}[a)]
			\item How high is the tension energy when jumping?
			\item What maximum height above the water does Lunerald reach?
			\item How much time does he have to do a somersault? (= flight duration)
		\end{enumerate}
		
		{\bf \item Estela is ready to bungee jump.} She stands on a $h_k=\SI{80}{\meter}$ high crane, the elastic rope is unstretched $l_s=\SI{30}{\meter }$ long. To be on the safe side, the fall $\SI{15}{\meter}$ should be stopped above the ground.
		\begin{enumerate}[a)]
			\item How long does the free fall last and what is its maximum speed?
			\item What height $h_r$ above the ground does it reach during the "rebound", i.e. when rising again through the rubber rope, when 40% of the energy is lost due to air resistance and deformation work?
			\item What must be the minimum spring hardness of the rope if Estela's mass is $m=\SI{60}{\kilogram}$?
		\end{enumerate}
		\fbox{
			\begin{minipage}{0.5\textwidth}
				For the solution, please \href{https://www.okuyakl.de/phys/p9frfL075/le075.pdf}{click here} or scan the QR code.\\
				Additional worksheets are available at
				
				\href{http://www.okuyakl.com}{www.okuyakl.com}
			\end{minipage}
			\hfill
			\begin{minipage}{0.4\textwidth}
				\includegraphics[width=1.5 cm]{../../viecher/zwe03}
				\includegraphics[width=3 cm]{qre075}
				\includegraphics[width=2 cm]{../../viecher/afanticon1}
				
		\end{minipage}}
		
	\end{enumerate}
	
\end{document}